\section{Saran}

Saran penelitian dibagi menjadi penerapan praktis dan pengembangan lanjutan.

Untuk penerapan praktis, CoDA-EDA sesuai digunakan ketika prioritas utama adalah efisiensi memori dan pembaruan kebijakan berkala pada perangkat \textit{edge}. 
Jika prioritasnya adalah kualitas penambangan statis pada dataset berdimensi besar, algoritma berpopulasi seperti ACO dan UBK dapat menjadi alternatif. 
Jumlah aturan pada strategi \textit{separate-and-conquer} juga perlu disesuaikan karena peningkatan anggaran aturan dapat memperluas \textit{coverage} dengan konsekuensi waktu komputasi dan kompleksitas kebijakan yang lebih tinggi.

Untuk pengembangan lanjutan, ambang \textit{drift} ($\tau$) dapat dibuat adaptif terhadap karakteristik distribusi agar tetap efektif pada dataset berdimensi besar. 
Selain itu, fungsi \textit{fitness} atau mekanisme pencarian dapat dikembangkan agar lebih sadar \textit{coverage}, sehingga kualitas tidak terlalu bergantung pada jumlah aturan. 
Pemodelan dependensi orde lebih tinggi secara selektif juga layak dieksplorasi dengan tetap mempertahankan efisiensi memori. 
Terakhir, validasi perlu diperluas ke lebih banyak dataset independen, data operasional IoT nyata, beragam perangkat \textit{edge}, serta integrasi dalam \textit{pipeline} pengelolaan kebijakan untuk menilai kelayakan penerapan secara lebih menyeluruh.